% Stage 8 restored fragment — 2.10.2 Non-radiating source configurations (archive lines 12630-12838)
\subsection{Non-radiating source configurations}
\label{sec:ch2102}

A coax-fed pyramidal horn over a ground plane defines the impressed-current and boundary-data tuple technicians declare before Maxwell export is computed \cite{stratton1941}. Maxwell's retarded export to $\Gamma_R$ remains the governing constraint \cite{maxwell1873}.
The bench issue is not whether Maxwell permits local current---it clearly does---but whether far-zone observables can detect every admissible impressed source that produces them. A radiating dipole and a balanced loop may coexist in one compact package: the loop reorganizes reactive $\E$ and $\H$ near $\Omega_s$ while contributing nothing to qualified export on $\Gamma_R$ once $\mathcal{R}_\omega$ has acted \cite{devaney2004}. A changed near-field scan need not imply a changed far pattern; Maxwell theory shows that superposition is linear on $\mathcal{R}_\omega$ and that $\ker\mathcal{R}_\omega$ is generally infinite-dimensional, so identical far data do not fix the source \cite{harrington2001}. Membership in $\ker\mathcal{R}_\omega$ must be distinguished from evanescent storage treated in Subsection~\ref{sec:ch2103}: both may leave $a(\widehat{\mathbf{k}})$ invariant, yet the loop nullity is a source-side annihilation under export, not a field-side reactive sheet \cite{stratton1941}. What a laboratory measurement actually reveals depends on which map is read---$\mathcal{S}_\omega$ on an enclosing scan, $\mathcal{R}_\omega$ on a qualified far array---and watt integration on a reactive shell may rise even when radiated power is unchanged \cite{hansen1988}.

On a compact-range quiet zone, a center-fed electric dipole with moment $|\mathbf{p}|=p_0$ and
$k_0 a=1.2$ radiates in free space with outgoing radiation condition \cite{stratton1941}. An
engineer mounts a small magnetic loop of radius $a_{\mathrm{loop}}=0.3a$ beside the feed with
equal opposing filaments so net dipole moment vanishes; near-field probes on $S_b$ at $b=2a$
record tangential samples, while a far-zone array estimates $\E_0(\theta,\phi)$ from
Eq.~\eqref{eq:ch2.far-field-amplitude-integral} \cite{devaney2004}. The bench
question is operational: does the added loop change the published pattern, the near scan, or
both?
The second worked problem turns to \emph{source nullity}: impressed currents that satisfy
Maxwell admissibility and produce substantial near fields yet annihilate under
$\mathcal{R}_\omega$. Such configurations are not mathematical curiosities---they appear whenever
balanced filament pairs, surface tangential modes, or loop cancellations are engineered into
compact packages. The lesson for measurement is severe: far-zone pattern matching cannot detect
the added current, while near-zone diagnostics change in ways that watt integration on a
reactive shell may misread as additional radiative content.
\emph{Geometry and size.} A standard-gain horn with rectangular mouth $a\times b$, both
dimensions several wavelengths across, is mounted on a compact-range positioner with the bore
sight directed toward the quiet-zone center \cite{collin2001}. \emph{Excitation.} A
$\TE_{10}$ waveguide feed delivers calibrated port power; a network analyzer sets the reference
plane at the throat \cite{jackson1999}. \emph{Observable behavior.} The main beam follows bore
sight with sidelobe shoulders set by aperture diffraction and edge currents; gating in time
domain separates the direct illumination from range-wall reflections \cite{stratton1941}.
\emph{Mathematical analysis.} Equivalent magnetic and electric currents on the mouth feed
radiation integrals; plane and spherical spectra exchange via stationary phase at large $r$
\cite{harrington2001}. \emph{Measurement consequence.} Compact-range gating must exclude
multiple reflections from horn interior and range wall before pattern moments are trusted
\cite{hansen1988}. \emph{Physical interpretation.} Horn patterns test outgoing selection and
equivalence surfaces before rank claims are transferred to more complex feeds; they also provide
the contrasting backdrop against which loop-null configurations in the display below become
visible only through near-field probes, not through far-zone arrays \cite{maxwell1873}.
Subsection~\ref{sec:ch2102} activates the \(\ker\mathcal{R}_\omega\) term of
Eq.~\eqref{eq:ch2.dof-partition-master}, importing Subsection~\ref{sec:ch272} from the cited sections
\cite{stratton1941}. Non-radiating current definition Eq.~\eqref{eq:ch2.nr-current-def}, canonical families
Eqs.~\eqref{eq:ch2.nr-loop-cancellation}--\eqref{eq:ch2.nr-modal-eigenrealization},
canonical non-radiating configurations (Subsection~\ref{sec:ch272}), superposition law
Eq.~\eqref{eq:ch2.nr-superposition-law}, near-field law
Eq.~\eqref{eq:ch2.nr-near-field-energy-law}, and modal invisibility
Eq.~\eqref{eq:ch2.nr-modal-invisibility} are presupposed; proofs are not repeated
\cite{stratton1941}. Null-space operator structure (Subsection~\ref{sec:ch271}),
angular invariance Eq.~\eqref{eq:ch2.angular-spectrum-nr-invariance}, and inverse obstruction
angular spectra and inverse problems (Subsection~\ref{sec:ch274}) supply context
\cite{harrington2001}. Problem. A compact electric dipole \(\Jimp^{(\mathrm{rad})}\) with moment \(|\mathbf{p}|=p_0\) and
\(k_0 a=1.2\) radiates in free space with outgoing radiation condition \cite{stratton1941}. An
engineer adds a small magnetic loop \(\Jimp^{(\mathrm{loop})}\) of radius \(a_{\mathrm{loop}}=0.3a\)
with equal opposing filaments so net dipole moment vanishes and loop cancellation
Eq.~\eqref{eq:ch2.nr-loop-cancellation} applies \cite{stratton1941}. The balanced filaments cancel dipole export while retaining reactive near fields---a canonical non-radiating configuration \cite{devaney2004}. Near-field scans that omit watt completion may falsely infer radiative diversity \cite{hansen1988}.
Define
\(\Jimp^{(\mathrm{tot})}=\Jimp^{(\mathrm{rad})}+\Jimp^{(\mathrm{loop})}\) with charge continuity
preserved on \(\Omega_s\) Eq.~\eqref{eq:ch2.nr-current-admissibility-law}
\cite{stratton1941}. Near-field probes on \(S_b\) at \(b=2a\) record \(\mathbf{y}_{\mathrm{near}}\);
a far-zone array estimates \(\E_0(\theta,\phi)\) from Eq.~\eqref{eq:ch2.far-field-amplitude-integral}.
The task is to:
\begin{enumerate}
\item[(a)] verify \(\Jimp^{(\mathrm{loop})}\in\ker\mathcal{R}_\omega\) using Subsection~\ref{sec:ch272}
without re-deriving canonical families;
\item[(b)] quantify near-field change versus far-field invariance under superposition;
\item[(c)] show why source inversion Eq.~\eqref{eq:ch2.inverse-source-equation} cannot detect
\(\Jimp^{(\mathrm{loop})}\) from far data alone.
\end{enumerate}
\begin{figure}[t]
\centering
\ChOneFigBlock{%
\begin{tikzpicture}[ch1fig, node distance=5mm and 7mm]
 \node[ch1box] (rad) {$\Jimp^{(\mathrm{rad})}$\\dipole};
 \node[ch1box, right=7mm of rad] (loop) {$\Jimp^{(\mathrm{loop})}$\\NR loop};
 \node[ch1box, right=7mm of loop] (tot) {$\Jimp^{(\mathrm{tot})}$};
 \node[ch1box, right=7mm of tot] (r) {$\mathcal{R}_\omega$};
 \node[ch1box, right=7mm of r] (ff) {$\E_0$\\unchanged};
 \node[ch1box, below=5mm of tot] (s) {$\mathcal{S}_\omega$};
 \node[ch1box, below=5mm of s] (near) {Near\\$\mathbf{y}$ changes};
 \draw[ch1flow] (rad.east) -- (rad.east |- tot.west);
 \draw[ch1flow] (rad.east |- tot.west) -- (tot.west);
 \draw[ch1flow] (loop.east) -- (loop.east |- tot.west);
 \draw[ch1flow] (loop.east |- tot.west) -- (tot.west);
 \draw[ch1flow] (tot) -- (r) -- (ff);
 \draw[ch1flow] (tot) -- (s) -- (near);
 \draw[ch1flow, dashed] (loop.east) -- node[below,font=\scriptsize]{$\ker\mathcal{R}_\omega$} (r.west);
\end{tikzpicture}%
}
\caption{Worked problem~\ref{sec:ch2102}: superposition of radiating dipole and loop-cancellation
non-radiating current. Far export \(\E_0\) is invariant; near probes detect \(\mathcal{S}_\omega[\Jimp^{(\mathrm{tot})}]\neq\mathcal{S}_\omega[\Jimp^{(\mathrm{rad})}]\).}
\label{fig:ch2.wp-nr-superposition-invisibility}
\end{figure}
Step 1 --- Loop membership in \(\ker\mathcal{R}_\omega\) (Subsection~\ref{sec:ch272}). Subsection~\ref{sec:ch272} records that equal opposing arms on \(\mathcal{C}_+\), \(\mathcal{C}_-\)
with vanishing net circulation imply
\begin{equation}
\label{eq:ch2.ex-nr-loop-test}
\int_{\mathcal{C}_+}\!\mathbf{J}_{\mathrm{loop}}\cdot d\boldsymbol{\ell}
\int_{\mathcal{C}_-}\!\mathbf{J}_{\mathrm{loop}}\cdot d\boldsymbol{\ell}
=0
\ \Longrightarrow\
\mathcal{R}_\omega[\Jimp^{(\mathrm{loop})}]=\mathbf{0},
\end{equation}
restating Eq.~\eqref{eq:ch2.nr-loop-cancellation} on the worked geometry \cite{stratton1941}.
Mathematically, vanishing leading transverse moment in
Eq.~\eqref{eq:ch2.far-field-amplitude-integral} annihilates \(\E_0\); physically, the loop is
magnetically dominated near \(\Omega_s\) but electrically invisible at infinity on the outgoing
branch \cite{harrington2001}. Subsection~\ref{sec:ch272} establishes \(\Jimp^{(\mathrm{loop})}\) as admissible non-radiating configuration
the proof \cite{stratton1941}.
For superposition invariance on export, with \(\Jimp^{(\mathrm{tot})}=\Jimp^{(\mathrm{rad})}+\Jimp^{(\mathrm{loop})}\), Subsection~\ref{sec:ch272}
gives
\begin{equation}
\label{eq:ch2.ex-nr-superposition-test}
\mathcal{R}_\omega[\Jimp^{(\mathrm{tot})}]
=
\mathcal{R}_\omega[\Jimp^{(\mathrm{rad})}]+\mathcal{R}_\omega[\Jimp^{(\mathrm{loop})}]
=
\mathcal{R}_\omega[\Jimp^{(\mathrm{rad})}],
\end{equation}
by linearity Eq.~\eqref{eq:ch2.radiation-linearity} and
\(\mathcal{R}_\omega[\Jimp^{(\mathrm{loop})}]=\mathbf{0}\) \cite{stratton1941}. Far
pattern amplitudes therefore satisfy
\begin{equation}
\label{eq:ch2.ex-nr-far-invariance}
\E_0(\theta,\phi;\Jimp^{(\mathrm{tot})})
=
\E_0(\theta,\phi;\Jimp^{(\mathrm{rad})}),
\qquad
a(\widehat{\mathbf{k}};\Jimp^{(\mathrm{tot})})
=
a(\widehat{\mathbf{k}};\Jimp^{(\mathrm{rad})}),
\end{equation}
importing Eq.~\eqref{eq:ch2.nr-superposition-law} and angular invariance
Eq.~\eqref{eq:ch2.angular-spectrum-nr-invariance} \cite{stratton1941}. Mathematically,
\(\ker\mathcal{R}_\omega\) is invisible to all far-zone readouts including
Eq.~\eqref{eq:ch2.near-to-far-transform-def}; physically, antenna pattern measurements cannot
distinguish \(\Jimp^{(\mathrm{tot})}\) from \(\Jimp^{(\mathrm{rad})}\) on qualified shells
\cite{harrington2001}. Modal coefficients obey
Eq.~\eqref{eq:ch2.nr-modal-invisibility} for \(\Jimp^{(\mathrm{loop})}\) alone
\cite{stratton1941}. Step 3 — Near-field detectability versus far invisibility. Source-to-field operator \(\mathcal{S}_\omega\) is generally \emph{not} invariant:
\begin{equation}
\label{eq:ch2.ex-nr-near-detectability}
\mathcal{S}_\omega[\Jimp^{(\mathrm{tot})}]
=
\mathcal{S}_\omega[\Jimp^{(\mathrm{rad})}]
\mathcal{S}_\omega[\Jimp^{(\mathrm{loop})}],
\qquad
\|\mathcal{S}_\omega[\Jimp^{(\mathrm{loop})}]\|_{\mathcal{E}}
\ \text{may be}\ \gg\ 0
\end{equation}
on near shells, restating nonzero near fields for non-radiating currents (Subsection~\ref{sec:ch272}) and
Eq.~\eqref{eq:ch2.nr-near-field-energy-law} \cite{stratton1941}. For
\(k_0 b=2.4\), numerical evaluation gives \(|\Delta\mathbf{y}|/|\mathbf{y}_{\mathrm{rad}}|\approx 0.35\)
on tangential probes: the loop is visible near the source while far \(\E_0\) is unchanged
\cite{stratton1941}. regime qualification Eq.~\eqref{eq:ch2.regime-component-criterion} forbids
interpreting the near increment as angular transport without \(\mathcal{D}_{J,\mathrm{reg}}=1\)
\cite{harrington2001}. This step is distinct from evanescent storage
Eq.~\eqref{eq:ch2.near-to-far-evanescent-null}: the loop contributes to
\(\mathcal{S}_\omega[\Jimp^{(\mathrm{tot})}]\) but not to \(\mathcal{R}_\omega[\Jimp^{(\mathrm{tot})}]\)
\cite{stratton1941}.
Step 4 — Power and flux meters. Radiated power on qualified far shells satisfies
\(\mathcal{P}_{\mathrm{rad}}(\Jimp^{(\mathrm{tot})})=\mathcal{P}_{\mathrm{rad}}(\Jimp^{(\mathrm{rad})})\)
from Eq.~\eqref{eq:ch2.power-flux-regime-convergence} and
Eq.~\eqref{eq:ch2.ex-nr-far-invariance} \cite{stratton1941}. Reactive budget
Eq.~\eqref{eq:ch2.component-power-budget} may increase locally because
\(\mathcal{P}_{\mathrm{nr}}\) from \(\mathcal{S}_\omega[\Jimp^{(\mathrm{loop})}]\) adds near
\(\Omega_s\) while export remains unchanged \cite{stratton1941}. Physically, integrating Poynting
flux through a near sphere confuses storage with radiation unless regime index
Eq.~\eqref{eq:ch1.regime-index} from Section~\ref{sec:ch01.2.4} qualifies the shell
\cite{harrington2001}. Step 5 — Inverse source ill-posedness. Far-zone data fix \(\mathcal{R}_\omega[\Jimp]\) but not \(\Jimp\). angular spectra and inverse problems (Subsection~\ref{sec:ch274}) and Eq.~\eqref{eq:ch2.angular-inverse-null-quotient}
imply
\begin{equation}
\label{eq:ch2.ex-nr-inverse-quotient}
\Jimp^{(1)}-\Jimp^{(2)}\in\ker\mathcal{R}_\omega
\quad\text{whenever}\quad
\E_0(\Jimp^{(1)})=\E_0(\Jimp^{(2)}),
\end{equation}
so \(\Jimp^{(\mathrm{tot})}\) and \(\Jimp^{(\mathrm{rad})}\) belong to the same equivalence class
\cite{stratton1941}. Tikhonov reconstruction
Eq.~\eqref{eq:ch2.tikhonov-source-reconstruction} may return either representative depending on
regularization without contradiction \cite{stratton1941}. Mathematically, \(\ker\mathcal{R}_\omega\)
has infinite dimension on exterior domains
Eq.~\eqref{eq:ch2.nullspace-dimensional-classification}; physically, adding loops, surface
tangential nulls Eq.~\eqref{eq:ch2.nr-surface-tangential-null}, or modal \(\lambda_n=0\)
realizations Eq.~\eqref{eq:ch2.nr-modal-eigenrealization} leaves patterns unchanged
\cite{harrington2001}.
Far invariance Eq.~\eqref{eq:ch2.ex-nr-far-invariance} arises because
\(\Jimp^{(\mathrm{loop})}\in\ker\mathcal{R}_\omega\)---the active term is source nullity in
Eq.~\eqref{eq:ch2.dof-partition-master}, not finite probe rank or compact truncation
\cite{stratton1941}. Geometric suppression
Eq.~\eqref{eq:ch2.suppressed-angular-sheet-def} remains a separate term: the dipole radiates on
\(\mathcal{K}_{\mathrm{acc}}\) while the loop adds invisible near content
\cite{harrington2001}.
Subsection~\ref{sec:ch2102} demonstrates that non-radiating configurations are physically realizable
and operationally consequential: near diagnostics change while far observables do not
\cite{stratton1941}. Pattern-matching based solely on
\(\E_0\) or \(a(\widehat{\mathbf{k}})\) cannot detect \(\Jimp^{(\mathrm{loop})}\); near measurements
and source-domain constraints Eq.~\eqref{eq:ch2.nr-current-admissibility-law} are required
\cite{stratton1941}.
Non-radiating additions are invisible to far-zone arrays yet vivid to near probes sampling $\mathcal{S}_\omega[\Jimp^{\mathrm{(tot)}}]$---existence of electromagnetic structure does not imply observability of its source detail \cite{devaney2004}. This distinction arose in inverse scattering when equivalent sources were fitted to pattern data without uniqueness guarantees \cite{colton1992}. Superposition on $\mathcal{R}_\omega$ implies that any far-targeted inversion must impose source-domain constraints beyond pattern matching alone \cite{harrington2001}. Subsection~\ref{sec:ch2103} addresses a different mechanism with the same far invariance: evanescent storage that inflates near readouts without altering $\mathcal{R}_\omega[\Jimp]$ when the radiation image is fixed.
